\bookchapter{第五十五章：項目二——天氣查詢 App（網絡版）}

\chapterintro{記賬本係「練內功」，天氣 App 係「見世面」——佢要同互聯網打交道。呢個項目你會學到：申請免費 API、解析 JSON、處理「載入中／成功／失敗」三種介面狀態，同埋點樣將 API 密匙收埋好。做完佢，你就具備咗整任何「上網 App」嘅能力——新聞、匯率、翻譯，套路一模一樣。}

\booksection{一、申請免費天氣 API}

個人練手，推薦 OpenWeatherMap 或者國內嘅和風天氣（QWeather）——都有免費額度，註冊即刻用得。以 OpenWeatherMap 為例：

\begin{enumerate}
  \item 註冊賬號 喺 API keys 頁面生成一個 Key（apikey）；
  \item 免費版可以調用 Current Weather（當前天氣）接口：api.openweathermap.org/data/2.5/weather? q=beijing\&appid=你嘅Key\&units=metric；
  \item 喺瀏覽器度先試一次，睇吓返回嘅 JSON 係咩樣——呢個係接任何 API 之前必做嘅「預習」。
\end{enumerate}

\textbf{避坑提醒：密匙絕對唔可以寫死喺程式碼入面}\par

API Key 等於你嘅「銀包密碼」，寫入程式碼一提交上 GitHub，幾個鐘之內就會被爬蟲掃走盜刷。正確做法：擺入 local.properties（呢個檔案已經被 gitignore 忽略），再透過 Gradle 嘅 buildConfig 讀入程式碼。搜「Android 隱藏 API Key 嘅最佳實踐」跟住改就得。

\booksection{二、網絡層：Retrofit + JSON 解析}

先照 API 返回嘅 JSON 定義「資料模型」，對唔上嘅欄位一律忽略：

\begin{lstlisting}[language=Kotlin]
// JSON: {"weather":[{"description":"晴"}],"main":{"temp":26.5,"humidity":40},"name":"北京"}
data class WeatherResponse(
    val weather: List<Weather>,
    val main: Main,
    val name: String
)
data class Weather(val description: String)
data class Main(val temp: Double, val humidity: Int)

interface WeatherApi {
    @GET("data/2.5/weather")
    suspend fun getWeather(
        @Query("q") city: String,
        @Query("appid") key: String,
        @Query("units") metric: String = "metric",
        @Query("lang") lang: String = "zh_cn"
    ): WeatherResponse
}
\end{lstlisting}

Gson 會自動將 JSON 塞入呢啲 data class——欄位名一致，魔法自動發生。

\booksection{三、三種介面狀態：載入中 / 成功 / 失敗}

上網 App 嘅靈魂在於「狀態管理」。網絡慢、斷網、城市名打錯……每種情況介面都要畀用戶一個交代。用一個 sealed class（密封類）優雅咁表示：

\begin{lstlisting}[language=Kotlin]
sealed class UiState {
    object Loading : UiState()                    // 兜緊圈
    data class Success(val data: WeatherResponse) : UiState()  // 成功
    data class Error(val message: String) : UiState()          // 失敗
}

// ViewModel 入面：
fun search(city: String) {
    viewModelScope.launch {
        _state.value = UiState.Loading
        _state.value = try {
            UiState.Success(api.getWeather(city))
        } catch (e: Exception) {
            UiState.Error("搵唔到呢個城市，轉個名再試吓？")
        }
    }
}

// 界面入面：
when (val s = state) {
    is UiState.Loading -> CircularProgressIndicator()   // 兜圈
    is UiState.Success -> WeatherContent(s.data)        // 展示
    is UiState.Error -> Text(s.message)                 // 報錯
}
\end{lstlisting}

呢個 Loading/Success/Error 三段式係上網介面嘅「萬金油模板」，先背低佢，第日所有上網 App 照住個模畫個樣。

\booksection{四、實戰介面程式碼：將天氣畫出嚟}

狀態有咗，介面就係咁樣（關鍵片段）：

\begin{lstlisting}[language=Kotlin]
@Composable
fun WeatherContent(data: WeatherResponse) {
    Column(
        Modifier.fillMaxSize().padding(24.dp),
        horizontalAlignment = Alignment.CenterHorizontally
    ) {
        Text(data.name, fontSize = 28.sp, fontWeight = FontWeight.Bold)
        Text(
            "${data.main.temp}°C",
            fontSize = 72.sp,
            fontWeight = FontWeight.Light
        )
        Text(
            data.weather.firstOrNull()?.description ?: "未知",
            fontSize = 20.sp,
            color = Color.Gray
        )
        Spacer(Modifier.height(16.dp))
        Row {
            Text(" 溼度 ${data.main.humidity}%")
            Spacer(Modifier.width(24.dp))
            Text(" 體感 ${data.main.temp}°C")
        }
    }
}
\end{lstlisting}

幾個小細節值得學：firstOrNull() 安全攞第一個元素（列表可能係空嘅），?: “未知” 畀個兜底文案——空安全原則喺介面入面同樣適用。溫度用 72.sp 嘅大字號撐場面，「輕」字重（Light）令個數字望落優雅，呢個係天氣類 App 嘅經典配方。

\booksection{五、介面同互動細節}

\booksection{七、介面同互動細節}

\begin{itemize}
  \item 搜尋框：頂部 TextField，輸入城市名撳搜尋（鍵盤上面設「搜尋」動作，一撳即搜）；
  \item 天氣卡：大字號溫度（26.5°）、天氣描述（晴）、濕度、城市名；
  \item 收藏城市：將成日睇嘅城市儲入 DataStore，下次打開自動幫你查——順手溫習返第 6 章；
  \item 下拉刷新：喺 Compose 用 PullToRefreshBox（Material 3 提供），手感即刻「專業」起上嚟。
\end{itemize}

\begin{tcbinfo}{小貼士：圖片令界面活起嚟}
幫每種天氣配一個 emoji 或者圖標（），用 when(description) 嚟對應——三行代碼，界面嘅溫度感即刻翻倍。亦都可以載入第三方嘅天氣圖標庫，但練手階段 emoji 已經夠晒用。

\end{tcbinfo}

\begin{tcbpractice}{嚟動手：天氣 App 驗收清單}
(1) 正常搜得出天氣；(2) 斷開網絡，App 優雅報錯唔會崩潰（開飛行模式嚟測！）；(3) 輸入一個唔存在嘅城市，提示都要友好；(4) 收藏兩個城市，重啟 App 之後仲喺度；(5) 轉屏唔會冇咗狀態。五條全部過關，呢個項目先算畢業。

\end{tcbpractice}

\begin{tcbsummary}{本章小結}
\begin{itemize}
  \item 先用瀏覽器試 API、睇明 JSON，先至寫代碼——「預習」永遠唔蝕底
  \item API Key 收埋入 local.properties，千祈唔好入代碼庫
  \item Retrofit 接口 + data class：欄位對得返，JSON 自動解析
  \item Loading/Success/Error 三段式狀態係聯網界面嘅萬金油
  \item 斷網、空輸入、城市唔存在——錯誤處理係聯網 App 嘅及格線
\end{itemize}

\end{tcbsummary}
